\documentclass[11pt]{article}
\usepackage[a4paper,margin=24mm]{geometry}
\usepackage[T1]{fontenc}
\usepackage{lmodern,amsmath,amssymb,amsthm,mathtools,microtype,enumitem}
\usepackage[colorlinks=true,linkcolor=blue,urlcolor=blue]{hyperref}
\hypersetup{pdftitle={The recurrence from one telescope},pdfauthor={Mathathon peer-review digest}}
\newcommand{\T}{\mathcal T}
\newcommand{\q}[1]{\medskip\noindent\textbf{#1}}
\setlist{itemsep=1pt,topsep=3pt}
\title{The recurrence from one telescope\\[2pt]
\large a shorter replacement for Section 2 of ``Why the maximum-principle proof of $\zeta(3)$ works''}
\author{Mathathon peer-review digest (Claude)}
\date{15 November 2026 \quad(draft)}
\begin{document}
\maketitle

\noindent Section~2 of the expanded exposition derives the recurrence in four pages: a second index
$m$, a telescope $G$ for the rows, a different telescope $g$ for the diagonal, and an elimination.
Here the same result comes from \emph{one} telescope used twice. There is no second index. The
output is exactly what Section~3 of that exposition needs (its equations (11) and (13)).

\q{Notation, as there.}
$r_n(t)=\dfrac{(1-t)\cdots(n-t)}{t(t+1)\cdots(t+n)}$,\quad $\T(f)=-\sum_{t\ge1}f'(t)$,\quad
$I_n=\T(r_n^2)$. The one analytic fact:
\begin{equation}\label{tel}
\T\bigl(g(t+1)-g(t)\bigr)=0
\end{equation}
if $g$ is a rational function with no pole at $1,2,3,\dots$, a zero of order at least two at
$t=1$, and a finite limit at infinity.

\q{What we want.} With the current $J_n=n^3(I_{n-1}-I_n)$ through the link from $n-1$ to $n$:
\begin{equation}\label{K}
J_n-J_{n+1}=4(2n+1)^3\,I_n\qquad(n\ge1).
\end{equation}

\section*{1. Neighbours share a factor}

From the products, $r_n=\frac{n-t}{n+t}\,r_{n-1}$. So the two functions at the ends of the link
from $n-1$ to $n$ are multiples of one function, which belongs to the link:
\[
r_{n-1}=(n+t)\,v_n,\qquad r_n=(n-t)\,v_n,\qquad
v_n(t)=\frac{(1-t)\cdots(n-1-t)}{t(t+1)\cdots(t+n)} .
\]
Adding and subtracting, $v_n=\frac{r_{n-1}+r_n}{2n}$: the link function is the mean of its two ends.
Besides the values $I_{n-1},I_n$ at the ends we shall need one number in the middle of the link,
\[
e_n=\T(r_n\,r_{n-1})=\T\bigl((n^2-t^2)\,v_n^2\bigr).
\]

\section*{2. The telescope}

To use \eqref{tel} we need functions with a finite limit at infinity. The link function $v_n$ decays
like $t^{-2}$; multiplying it by $n^2-t^2$ restores the balance and gives
\[
s_n(t)=(n^2-t^2)\,v_n(t)=\frac{(1-t)(2-t)\cdots(n-t)}{t(t+1)\cdots(t+n-1)} ,
\]
with $n$ zeros, at $1,\dots,n$, over $n$ poles, at $0,\dots,-(n-1)$. Now shift $t$ by one. The zeros
move to $0,\dots,n-1$ and the poles to $-1,\dots,-n$:
\[
s_n(t+1)=\frac{(-t)(1-t)\cdots(n-1-t)}{(t+1)\cdots(t+n)}=-t^2\,v_n(t).
\]
\emph{So $s_n$ and its shift are both multiples of the same $v_n$:} the factors are $n^2-t^2$ and
$-t^2$. Hence for any two links $n,m\ge1$ the product $g=s_ns_m$ satisfies
\[
g(t+1)-g(t)=\bigl[t^4-(n^2-t^2)(m^2-t^2)\bigr]v_nv_m=\bigl[(n^2+m^2)\,t^2-n^2m^2\bigr]\,v_nv_m .
\]
This $g$ has a double zero at $t=1$ (one from each factor), no pole at a positive integer, and the
limit $\pm1$. By \eqref{tel},
\begin{equation}\label{star}
\boxed{\;\T\Bigl(\bigl[(n^2+m^2)\,t^2-n^2m^2\bigr]\,v_n\,v_m\Bigr)=0\qquad(n,m\ge1).\;}
\end{equation}
This is the only telescope in the derivation. We use it for a link with itself, and for two links
that meet at a node.

\section*{3. A link with itself ($m=n$): the middle of the link}

\eqref{star} reads $\T\bigl((2t^2-n^2)\,v_n^2\bigr)=0$. On the other hand, by Section~1,
\[
r_{n-1}^2+r_n^2-6\,r_nr_{n-1}=\bigl[2(n^2+t^2)-6(n^2-t^2)\bigr]v_n^2=4\,(2t^2-n^2)\,v_n^2 .
\]
Therefore
\begin{equation}\label{six}
I_{n-1}+I_n=6\,e_n :
\end{equation}
the value in the middle of a link is one sixth of the sum of the values at its ends.

\section*{4. Two links at a node ($m=n+1$): the node}

The links $n$ and $n+1$ meet at the node $n$, and $r_n=(n-t)\,v_n=(n+1+t)\,v_{n+1}$. So
$v_nv_{n+1}=\dfrac{r_n^2}{(n-t)(n+1+t)}$, and the bracket in \eqref{star}, divided by these two
factors, splits into partial fractions. The numerator $(n^2+(n+1)^2)\,t^2-n^2(n+1)^2$ takes the
values $n^4$ and $(n+1)^4$ at the two poles $t=n$ and $t=-n-1$, so
\[
\frac{(n^2+(n+1)^2)\,t^2-n^2(n+1)^2}{(n-t)(n+1+t)}
=-\bigl(n^2+(n+1)^2\bigr)+\frac1{2n+1}\Bigl[\frac{n^4}{n-t}+\frac{(n+1)^4}{n+1+t}\Bigr].
\]
The two simple fractions give back the link functions: $\frac{r_n^2}{n-t}=r_nv_n=\frac{r_n(r_n+r_{n-1})}{2n}$
and $\frac{r_n^2}{n+1+t}=r_nv_{n+1}=\frac{r_n(r_{n+1}+r_n)}{2(n+1)}$. Applying $\T$, \eqref{star}
becomes
\begin{equation}\label{node}
n^3\,(e_n+I_n)+(n+1)^3\,(e_{n+1}+I_n)=2\,(2n+1)\,\bigl(n^2+(n+1)^2\bigr)\,I_n .
\end{equation}
\emph{This is where the cubes come from:} a fourth power at each pole, divided by the $2n$ or
$2(n+1)$ of the mean.

\section*{5. Removing the middles}

By \eqref{six}, $e_n+I_n=\frac16(I_{n-1}-I_n)+\frac43I_n$ and
$e_{n+1}+I_n=-\frac16(I_n-I_{n+1})+\frac43I_n$. In terms of the currents,
\[
n^3(e_n+I_n)=\tfrac16J_n+\tfrac43n^3I_n,\qquad (n+1)^3(e_{n+1}+I_n)=-\tfrac16J_{n+1}+\tfrac43(n+1)^3I_n .
\]
Put this into \eqref{node} and multiply by $6$. With $p=n+1$ and $q=n$,
\[
J_n-J_{n+1}=\bigl[12\,(p+q)(p^2+q^2)-8\,(p^3+q^3)\bigr]I_n=4\,(p+q)\bigl[3(p^2+q^2)-2(p^2-pq+q^2)\bigr]I_n
=4\,(p+q)^3\,I_n ,
\]
which is \eqref{K}. In the notation of the exposition, with $F_n=-J_n=n^3(I_n-I_{n-1})$, this is
$F_{n+1}=F_n+4(2n+1)^3I_n$, its equation (11).

\section*{What changed, and what is lost}

\begin{itemize}
\item \emph{One telescope instead of two, and no second index.} Both the $6$ and the cubes come out
of \eqref{star}: the $6$ from a link paired with itself, the cubes from two links paired at a node.
\item \emph{No design step.} In the exposition one has to decide which factors a polynomial
prefactor should cancel, twice. Here the function $s_n$ is forced: it is the multiple of the link
function that has a finite limit.
\item \emph{One partial fraction, with simple poles,} as before.
\item \emph{Lost:} the table $I_{n,m}$ and with it the remark that its row $m=0$ is the failed
approximation by partial sums. Section~3.3 of the exposition recovers that remark in another way
(``what happens if we remove the leakage'').
\item \emph{A picture that comes for free.} Written with differences, \eqref{six} and \eqref{node}
are
\[
(I_{n-1}-e_n)-(e_n-I_n)=4\,e_n,\qquad n^3(e_n-I_n)-(n+1)^3(I_n-e_{n+1})=2n(n+1)(2n+1)\,I_n .
\]
So the staircase $I_0,e_1,I_1,e_2,I_2,\dots$ is itself a chain of the kind studied in Section~3:
each link is split into two halves of conductance $n^3$, with a positive leak at every middle and
at every node. Section~5 above removes the middles from this finer chain.
\end{itemize}

\noindent\emph{Status.} Checked by the author only: the shift formula for $s_n$ symbolically for
$n\le5$; \eqref{star} in exact arithmetic for $1\le n,m\le7$; \eqref{six}, \eqref{node} and the
recurrence for $n\le12$; the partial fraction and the last line with $n$ (or $p,q$) symbolic. Three
AI agents are examining this and the other sections independently. The identity \eqref{star} is
equivalent to a relation between four neighbouring entries of the table $I_{n,m}$ found earlier in
this project by another agent; relations of this kind between neighbouring hypergeometric sums are
classical, and nothing here is claimed as new.
\end{document}
